\documentclass[11pt,a4paper]{article}

\usepackage[a4paper,top=24mm,bottom=25mm,left=26mm,right=24mm]{geometry}
\usepackage{fontspec}
\usepackage{amsmath,amssymb,bm}
\usepackage{booktabs}
\usepackage{tabularx}
\usepackage{array}
\usepackage{enumitem}
\usepackage{xcolor}
\usepackage{hyperref}
\usepackage{caption}

\setmainfont{Times New Roman}
\setsansfont{Helvetica Neue}
\setmonofont{Menlo}[Scale=0.84]
\definecolor{navy}{HTML}{17365D}
\definecolor{blue}{HTML}{1F5A94}
\definecolor{lightblue}{HTML}{EAF2F8}
\definecolor{darkgray}{HTML}{454545}
\hypersetup{
  colorlinks=true,
  linkcolor=navy,
  citecolor=navy,
  urlcolor=blue,
  pdftitle={Thiết kế và kiểm chứng PA-MPPI trên companion computer với ArduPilot},
  pdfauthor={Nguyễn Thành Lâm}
}

\pagestyle{plain}
\renewcommand{\abstractname}{Tóm tắt}
\renewcommand{\contentsname}{Mục lục}
\renewcommand{\tablename}{Bảng}
\renewcommand{\refname}{Tài liệu tham khảo}
\setlength{\emergencystretch}{2em}
\setlength{\parindent}{1.2em}
\setlength{\parskip}{0.25em}
\setlist{nosep,leftmargin=1.6em}
\captionsetup{font=small,labelfont=bf}

\newcommand{\code}[1]{\texttt{#1}}
\newcommand{\R}{\mathbb{R}}
\newcommand{\vect}[1]{\bm{#1}}
\newcolumntype{Y}{>{\raggedright\arraybackslash}X}
\newcolumntype{L}[1]{>{\raggedright\arraybackslash}p{#1}}

\title{\vspace{-1.5em}\textbf{Thiết kế và kiểm chứng PA-MPPI trên companion computer\vspace{0.2em}\\
với ArduPilot và Gazebo}}
\author{Nguyễn Thành Lâm\\[0.3em]
\small Báo cáo tiến độ kỹ thuật}
\date{13 tháng 09 năm 2026 -- evidence-first revision}

\begin{document}

\maketitle
\thispagestyle{empty}

\begin{abstract}
Báo cáo trình bày quá trình chuyển tầng tránh vật cản từ
\code{AP\_Avoidance} của ArduPilot sang companion computer. Hệ thống nhận
point cloud LiDAR 3D, tối ưu quỹ đạo bằng Model Predictive Path Integral
(MPPI), sau đó truyền setpoint cho ArduPilot ở chế độ \code{GUIDED}. Ba mức
triển khai được đánh giá: vanilla MPPI với lệnh vận tốc, PA-MPPI v0 với bản đồ
occupancy ba trạng thái và perception cost, và nhánh rigid-body experimental
với trạng thái $[\vect p,\vect q,\vect v,\vect\omega]$ cùng lệnh collective
thrust/body rate. Bản velocity-level có time-indexed global reference đã đạt
terminal trong ba lần chạy Gazebo/SITL (right-angle, narrow-gate và slalom,
mỗi map một seed). Nhánh rigid-body đã vượt qua kiểm tra hover, quaternion,
packet MAVLink và closed-loop clear-air offline, nhưng chưa qua hover gate live
trong Gazebo. Kết quả hỗ trợ tính khả thi của kiến trúc, không phải chứng nhận
robustness hay an toàn phần cứng.
\end{abstract}

\noindent\textbf{Từ khóa:} MPPI, PA-MPPI, quadrotor, ArduPilot, MAVLink,
LiDAR 3D, Gazebo, companion computer.

\vspace{0.8em}
\noindent\fcolorbox{navy}{lightblue}{\parbox{0.92\textwidth}{
\textbf{Trạng thái bằng chứng tại ngày 13/09/2026:}
\textbf{VERIFIED OFFLINE}: 40/40 unit tests và benchmark analytical 3 map
$\times$ 3 seed. \textbf{VERIFIED GAZEBO/SITL, ONE SEED}: velocity-level
right-angle, narrow-gate và slalom. \textbf{IMPLEMENTED, NOT LIVE-VERIFIED}:
rigid body-rate/thrust và PA-MPPI unknown-space ablation.
\textbf{PENDING GAZEBO}: rigid hover/body-rate và C-wall ablation.
\textbf{PENDING EDGE/HARDWARE}: runtime/thermal, sensor model, HIL và flight.}}

\clearpage
\begingroup
\small
\tableofcontents
\endgroup
\clearpage

\section{Bối cảnh và mục tiêu}

Pipeline tuần trước đưa obstacle vào \code{AP\_Avoidance} thông qua MAVLink.
Cách này yêu cầu companion computer chuyển dữ liệu cảm biến sang representation
và vòng đời mà ArduPilot quản lý. Tuy nhiên, protocol không bắt buộc tạo
persistent obstacle ID: \code{OBSTACLE\_DISTANCE\_3D} cho phép
\code{UINT16\_MAX} khi ID không biết \cite{mavlinkArdupilotmega2026}; target
handler cũng không đọc field này \cite{ardupilotSourceF808}. Mentor đề xuất tách
hoàn toàn local planner khỏi autopilot: companion computer tự xử lý perception,
sinh trajectory và gửi lệnh điều khiển; ArduPilot chỉ đảm nhiệm ổn định bay và
tracking setpoint.

Mục tiêu kỹ thuật của giai đoạn này gồm:

\begin{enumerate}
  \item xây dựng baseline MPPI không phụ thuộc \code{AP\_Avoidance};
  \item tích hợp LiDAR 3D, state estimation và MAVLink thành một vòng kín;
  \item bổ sung perception-aware objective theo PA-MPPI \cite{zhai2026pamppi};
  \item tiến dần từ velocity-level model tới rigid-body model;
  \item định nghĩa tiêu chí kiểm thử có thể tái lập trước khi chạy trên phần cứng.
\end{enumerate}

Phạm vi hiện tại là Gazebo và ArduPilot SITL. Báo cáo không xem kết quả offline
hoặc SITL là bằng chứng đủ cho flight test trên quadrotor thật.

\subsection*{Định danh phiên bản và quy tắc nguồn}

\begingroup\small
Project được audit tại commit
\texttt{910c6a824fad0ca38615f93a0a17fc2c36d591cf}. ArduPilot là dependency ngoài
repository, không phải vendored/submodule, tại commit
\texttt{f808f78ce5a518ca96f2fb36420608b2b6254367}; MAVLink gitlink tương ứng là
\texttt{71d925850d7ddb01d0e56defc1046f9d9d3bf8ae}. Protocol specification chỉ là
thẩm quyền về message fields; source đúng commit là thẩm quyền cuối về hành vi
firmware \cite{mavlinkCommon2026,ardupilotSourceF808}. Claim audit chi tiết nằm
trong \code{docs/SOURCE\_AUDIT.md}.
\endgroup

\section{Kiến trúc hệ thống}

\subsection{Phân tầng điều khiển}

Luồng dữ liệu triển khai trong repository được tóm tắt như sau:

\begin{center}
\fcolorbox{navy}{lightblue}{\parbox{0.88\textwidth}{\centering
LiDAR 3D + Gazebo odometry hoặc MAVLink telemetry\\
$\Downarrow$\\
lọc, downsample, đổi frame, occupancy mapping\\
$\Downarrow$\\
MPPI / PA-MPPI trên companion computer\\
$\Downarrow$\\
velocity target hoặc thrust/body-rate target qua MAVLink\\
$\Downarrow$\\
ArduPilot GUIDED $\rightarrow$ motor controller $\rightarrow$ Gazebo physics}}
\end{center}

Avoidance nội bộ của ArduPilot được tắt bằng \code{OA\_TYPE=0},
\code{AVOID\_ENABLE=0} và \code{PRX1\_TYPE=0}. Do đó point cloud không đi qua
\code{OBSTACLE\_DISTANCE\_3D}, \code{AP\_Proximity} hoặc obstacle database của
ArduPilot.

\subsection{Tại sao dùng kiến trúc này?}

ArduPilot cung cấp Guided setpoint interfaces và stabilization nội bộ
\cite{ardupilotGuidedDocs2026}; đồng thời có proximity/Simple Avoidance và
BendyRuler \cite{ardupilotSimpleAvoidance2026,ardupilotBendyRuler2026}.
\textbf{PROJECT DESIGN DECISION}: perception, occupancy mapping và PA-MPPI nằm
ở companion vì optimizer cần truy cập trực tiếp map ba trạng thái và perception
objective \cite{zhai2026pamppi}. Việc tắt AP OA trong experiment tránh hai
avoidance loop cùng sửa lệnh; nó không có nghĩa AP avoidance vốn không dùng được.

\begin{table}[ht]
\centering\small
\caption{Nguồn ràng buộc so với lựa chọn engineering.}
\begin{tabularx}{\textwidth}{L{0.21\textwidth}L{0.30\textwidth}Y}
\toprule
\textbf{Decision} & \textbf{Source-supported constraint} & \textbf{Our design choice / why}\\
\midrule
Companion planning & Guided nhận external setpoints; PA-MPPI dùng online map. & Giữ map/optimizer ngoài AP để không dịch internal representation.\\
Disable AP OA & AP OA là một stack độc lập. & Tắt trong ablation để tránh hai avoidance loop.\\
Velocity baseline & Guided hỗ trợ velocity/yaw-rate. & Reduced model để tách planner khỏi low-level validation.\\
Three-state map & Paper dùng unknown/free/occupied. & Fixed NumPy/Torch prototype; không gọi là ROG-Map.\\
Perception/unknown & Paper reward frontier và phạt non-free collision. & Port endpoint ray/LOS; giữ unknown collision policy.\\
Rigid branch & Paper dùng collective thrust/body rates. & Reduced closed-loop model, experimental only.\\
MAVLink interfaces & Spec/source hỗ trợ local velocity và target source nhận full rate vector. & Mask 1479 cho baseline; mask 128 + GUID=8 gate cho rigid.\\
50 Hz/rollouts & Paper: 50 Hz, 17,500 rollouts. & Rigid target 50 Hz, 512 rollouts do compute budget; phải benchmark live.\\
Safety gates & Low-level command giảm margin bảo vệ. & GUID readback, stale/brake/timeout, NaN và saturation gates.\\
\bottomrule
\end{tabularx}
\end{table}

\subsection{Các module phần mềm}

\begin{table}[ht]
\centering
\caption{Cấu trúc phần mềm chính.}
\begin{tabularx}{\textwidth}{L{0.34\textwidth}Y}
\toprule
\textbf{Module} & \textbf{Trách nhiệm} \\
\midrule
\code{mppi\_controller.py} & Vanilla MPPI, state 7 chiều và velocity command. \\
\code{occupancy\_grid.py} & Voxel map unknown/free/occupied và ray integration. \\
\code{pa\_mppi\_controller.py} & Perception-aware cost trên velocity-level model. \\
\code{rigid\_body\_pa\_mppi.py} & Rigid-body rollout, quaternion và thrust/body-rate command. \\
\code{mavlink\_interface.py} & Telemetry, velocity target và \code{SET\_ATTITUDE\_TARGET}. \\
\code{mppi\_local\_planner\_node.py} & Vòng điều khiển, waypoint, safety gate, diagnostics và RViz. \\
\bottomrule
\end{tabularx}
\end{table}

\section{Vanilla MPPI velocity-level}

\textbf{PROJECT DESIGN DECISION}: đây là tầng A trong staged validation. Nó
kiểm tra local planning trong khi vẫn giao attitude và motor stabilization cho
ArduPilot. Bằng chứng bên ngoài chỉ hỗ trợ việc Guided nhận velocity targets
\cite{ardupilotGuidedDocs2026,mavlinkCommon2026}; mô hình giảm bậc này không
được quy cho paper PA-MPPI.

\subsection{Mô hình dự đoán}

Baseline giữ ArduPilot ở tầng velocity tracking. State và control là

\begin{equation}
\vect x = [\vect p^\top,\vect v^\top,\psi]^\top \in \R^7,
\qquad
\vect u = [\vect v_{cmd}^\top,\dot\psi_{cmd}]^\top \in \R^4.
\end{equation}

Mô hình point-mass với first-order velocity lag được rời rạc hóa thành

\begin{align}
\vect v_{k+1} &= \vect v_k + \alpha(\vect v_{cmd,k}-\vect v_k),
& \alpha &= \min(\Delta t/\tau,1),\\
\vect p_{k+1} &= \vect p_k + \Delta t\,\vect v_{k+1},\\
\psi_{k+1} &= \psi_k + \Delta t\,\dot\psi_{cmd,k}.
\end{align}

Đây là mô hình giảm bậc có chủ ý. Attitude, motor dynamics và lực khí động được
đóng bên trong ArduPilot và Gazebo thay vì tối ưu trực tiếp. State rollout được
mở rộng thêm setpoint khả thi trước đó $u^{applied}_{k-1}$; low-pass, slew-rate
limit và saturation giống tầng MAVLink được áp dụng bên trong dynamics trước
khi tính $\vect v_{k+1}$. Đây là PROJECT DESIGN DECISION để prediction không
giả định bước nhảy control mà interface thực không gửi được.

\subsection{Tối ưu MPPI}

Tại mỗi chu kỳ, $N$ control sequence được tạo bằng cách thêm nhiễu Gaussian
vào nominal sequence. Với tổng cost $S_j$ của rollout thứ $j$, trọng số là

\begin{equation}
w_j = \frac{\exp(-(S_j-S_{\min})/\lambda)}
{\sum_{i=1}^{N}\exp(-(S_i-S_{\min})/\lambda)}.
\end{equation}

Control sequence mới là trung bình có trọng số của các rollout. Cost baseline
gồm khoảng cách tới waypoint, obstacle clearance, control effort, độ thay đổi
control, yaw alignment và terminal distance. Obstacle cost dùng softplus quanh
safe margin để giữ gradient chi phí mềm nhưng vẫn xử lý được point cloud rời rạc.
Sampling, importance weighting và receding-horizon update dựa trên MPPI/
information-theoretic MPC \cite{williams2018itmppi}; project không claim tái tạo
mọi theoretical correction term của implementation trong paper.

\subsection{Cost bám global path (bản velocity-level)}

Minařík et al. xây dựng cost reference-tracking bằng tổng cost input và độ đổi
input (Eq.~16), sau đó cộng metric giữa state rollout và reference state theo
thời gian (Eq.~17--18) \cite{minarik2024agilemppi}. Metric vị trí của paper là
weighted squared error; ngoài vị trí còn có quaternion, vận tốc và body-rate.
Đây là nguồn tham khảo cho việc bổ sung tracking cost, không phải bằng chứng
rằng mọi thành phần đó đã có trong baseline này.

Global planner/mission cung cấp polyline $\mathcal P$ trong ENU. Mỗi chu kỳ,
vị trí hiện tại được chiếu lên path để thu progress $s_0$; progress không được
lùi so với chu kỳ trước. Với tốc độ reference $v_r$, từng bước horizon dùng
$s_k=s_0+(k+1)v_r\Delta t$, từ đó nội suy $\vect p_k^{ref}$ và lấy sai phân hữu
hạn cho $\vect v_k^{ref}$. Profile paper-only dùng

\begin{equation}
  J_{\mathrm{ref}} = \sum_{k=0}^{H-1}\left[
  w_{p}
  \left(\frac{\|\vect p_k-\vect p_k^{ref}\|_2}{s_{\mathrm{path}}}\right)^2
  + w_v\|\vect v_k-\vect v_k^{ref}\|_2^2\right].
\end{equation}

Baseline project vẫn có tùy chọn nearest-polyline

\begin{equation}
  J_{\mathrm{polyline}} = w_{\mathrm{path}}
  \sum_{k=0}^{H-1}
  \left(\frac{d(\vect p_k,\mathcal P)}{s_{\mathrm{path}}}\right)^2,
  \qquad
  d(\vect p,\mathcal P)=\min_{\text{segment }e\in\mathcal P}
  \|\vect p-\operatorname{proj}_{e}(\vect p)\|_2 .
\end{equation}

với $d$ là projection lên segment gần nhất. Path được truyền qua
\code{--global-path}; \code{--goal} vẫn quyết định active waypoint và điều kiện
kết thúc.

Đây là \textbf{PROJECT DESIGN DECISION}, không phải reproduction đầy đủ của
Eq.~17--18: position/velocity đã time-indexed nhưng chưa có quaternion/body-rate
reference, dynamics rigid-body hay motor feasibility của paper. Polyline cũng
không phải collision certificate; path xuyên obstacle tạo xung đột tracking--
safety và phải được kiểm tra clearance trước khi chạy.

\subsection{Đối chiếu toàn bộ cost với paper}

Bảng dưới đây tách các thành phần có cơ sở từ Minařík et al. khỏi heuristic
được thêm cho bài toán velocity-level. Paper Eq.~16--18 và Eq.~22 là nguồn
tham khảo trực tiếp \cite{minarik2024agilemppi}; profile thử nghiệm
\code{mppi\_paper\_cost\_only.yaml} chỉ bật các hàng được đánh dấu.

\begin{table}[ht]
\centering\small
\caption{Đối chiếu cost project và profile paper-only.}
\label{tab:cost-audit}
\begin{tabularx}{\textwidth}{L{0.23\textwidth}L{0.31\textwidth}Y}
\toprule
\textbf{Cost hiện tại} & \textbf{Đối chiếu paper} & \textbf{Xử lý trong paper-only} \\
\midrule
Running goal distance & Paper dùng metric giữa state và reference state theo thời gian; không có active-waypoint heuristic riêng. & Tắt $w_{goal}$. \\
Terminal goal distance & Profile đã phạt reference state ở mọi bước; không cần active-waypoint terminal heuristic riêng. & Tắt $w_{terminal}$. \\
Softplus proximity & Không phải collision indicator Eq.~22. & Tắt; thay bằng $c_{obs}\mathbf{1}_{x\in\mathcal C_{obs}}$. \\
Control effort & Cùng nhóm với Eq.~16, nhưng paper dùng ma trận chéo $R$. & Dùng $R=\operatorname{diag}(0.01,0.05,0.05,0.10)$ sau ánh xạ velocity-level. \\
$R_{\Delta}$ input-change & Eq.~16 phạt $\Delta u_j^\mathsf{T}R_{\Delta}\Delta u_j$. & Dùng $R_{\Delta}=\operatorname{diag}(0.05,0.10,0.10,0.30)$ trên chuỗi action. \\
$w_{du}(u-v)$ & Không phải $\Delta u_j=u_{j+1}-u_j$ của Eq.~16. & Tắt; thay bằng $R_{\Delta}$ tính trên successive controls trong terminal callback. \\
Yaw heuristic & Không có term tương ứng trực tiếp trong Eq.~16--18. & Tắt. \\
Time-indexed global path & Paper dùng full reference state theo thời gian. & Bật $\vect p_k^{ref}$ và $\vect v_k^{ref}$; thiếu attitude/body-rate. \\
\bottomrule
\end{tabularx}
\end{table}

Trong profile paper-only, phần input-change được tính theo
\begin{equation}
J_{\Delta u}=\sum_{j=0}^{H-2}
\Delta u_j^\mathsf{T}R_{\Delta}\Delta u_j,
\qquad \Delta u_j=u_{j+1}-u_j,
\end{equation}
trên toàn bộ chuỗi control khả thi của mỗi rollout tại terminal callback. Vì action của
baseline là $[v_x,v_y,v_z,\dot\psi]$ thay vì $[F_t,\omega_d]$, ma trận
$R_{\Delta}$ của paper chỉ được ánh xạ theo thứ tự control này; đây là
adaptation có chủ đích, không phải khẳng định tương đương về vật lý.

\section{PA-MPPI v0}

\subsection{Occupancy map ba trạng thái}

PA-MPPI cần phân biệt vùng đã biết và chưa biết. Bản đồ voxel dùng quy ước

\begin{equation}
G(\vect p) \in \{-1,0,1\},
\quad
-1:\text{ unknown},\;0:\text{ free},\;1:\text{ occupied}.
\end{equation}
Đây là ba trạng thái và unknown-space safety policy từ paper
\cite{zhai2026pamppi}; backend dưới đây là implementation riêng của project.

Với mỗi LiDAR endpoint, voxel trên ray được đánh dấu free và endpoint được đánh
dấu occupied. Occupied được ghi sau để không bị ray lân cận xóa. Vì point cloud
Gazebo sau downsample chỉ giữ hit endpoint, mapper hiện đánh dấu một vùng gần
sensor là free. Đây là giả thiết thực dụng cho prototype, không thay thế việc
tích hợp đầy đủ max-range/miss ray của sensor thật.

\subsection{Perception-aware objective}

Khi chưa có line-of-sight trực tiếp tới goal, terminal cost được bổ sung bằng

\begin{equation}
\ell_{perception} = c_{PoI}(1-\langle \hat{\vect x}_B,
\hat{\vect d}_{goal}\rangle)^2
+ c_{occ}\,\bm{1}_{G(\vect r(t^*))=1}
+ c_{unk}\,\bm{1}_{G(\vect r(t^*))=-1}.
\end{equation}

Thành phần đầu hướng trục camera/body-x về goal. Hai thành phần sau phạt ray
đụng occupied voxel và thưởng endpoint nhìn vào unknown frontier. Collision
cost coi mọi voxel không phải free là không an toàn. Perception cost được tắt
khi goal đã nhìn thấy trực tiếp, phù hợp cấu trúc hai pha trong
\cite{zhai2026pamppi}.

PA-MPPI v0 vẫn dùng velocity-level dynamics. Mục đích của v0 là kiểm chứng
occupancy/perception integration trên interface ổn định trước khi hạ xuống tầng
thrust/body rate.
Đây là tầng B và là \textbf{PROJECT DESIGN DECISION}: unit test xác nhận v0 dùng
đúng cùng baseline dynamics, nhằm tách câu hỏi perception-aware planning khỏi
câu hỏi low-level thrust/body-rate.

\section{Rigid-body PA-MPPI experimental}

Đây là tầng C. Control space tiến gần formulation của paper hơn
\cite{zhai2026pamppi}, nhưng prediction model, mapper, rollout scale và validation
không khớp hoàn toàn; vì vậy tên đúng là ``experimental'', không phải full
reproduction.

\subsection{State, control và hệ trục}

Nhánh experimental dùng

\begin{equation}
\vect x=[\vect p_W^\top,\vect q_{WB}^\top,\vect v_W^\top,
\vect\omega_B^\top]^\top\in\R^{13},
\quad
\vect u=[c,\vect\omega_{cmd}^\top]^\top\in\R^4.
\end{equation}

Quaternion $\vect q_{WB}=[q_w,q_x,q_y,q_z]$ biến đổi vector body FLU của
Gazebo sang world ENU. Control $c$ là collective thrust theo newton;
$\vect\omega_{cmd}$ là body rate FLU. Trước khi gửi MAVLink, phép đổi

\begin{equation}
[p,q,r]_{FRD}=[p,-q,-r]_{FLU}
\end{equation}

đưa body rate về quy ước FRD của ArduPilot.

\subsection{Dynamics rời rạc}

Với $\vect e_3=[0,0,1]^\top$, gia tốc trong ENU là

\begin{equation}
\dot{\vect v}_W = \frac{c}{m}R(\vect q_{WB})\vect e_3
+[0,0,-g]^\top.
\end{equation}

Quaternion kinematics và body-rate tracking được mô hình hóa bởi

\begin{align}
\dot{\vect q}_{WB} &= \frac{1}{2}\vect q_{WB}\otimes[0,\vect\omega_B],\\
\vect\omega_{k+1} &= \vect\omega_k+
\alpha_\omega(\vect\omega_{cmd,k}-\vect\omega_k),
&\alpha_\omega&=\min(\Delta t/\tau_\omega,1).
\end{align}

Forward Euler được dùng cho velocity, quaternion và rate; position dùng thêm
$\tfrac12\vect a\Delta t^2$. Quaternion được chuẩn hóa sau mỗi bước. Model hiện
mô tả response của closed-loop body-rate controller bằng first-order lag. Phần
vật lý chỉ gồm thrust acceleration và quaternion kinematics; phần inner-loop là
approximation của ArduPilot. Nó chưa mô phỏng torque, moment of inertia, rotor
allocation, drag, battery hoặc motor lag tách biệt như dynamics trong paper
\cite{zhai2026pamppi}. Đây là reduced closed-loop prediction model; không có
inertia value nào được suy đoán.

Khối lượng 2.10 kg được cộng từ SDF: iris/IMU/bốn rotor
$1.5+0.15+4(0.025)=1.75$ kg, gimbal 0.23 kg và sensor pod 0.12 kg. Giá trị phụ
thuộc đúng include tree; center of mass và inertia tổng chưa được nhận dạng.
Hover force là $c_h=mg$. Ánh xạ từ newton sang thrust
chuẩn hóa dùng

\begin{equation}
T_{norm}=T_{hover,norm}\frac{c}{mg},
\end{equation}

với $T_{hover,norm}=0.38$. Đây là calibration assumption từ historical log
\code{logs/00000002.BIN}, trong đó ArduPilot từng học
\code{MOT\_THST\_HOVER} khoảng 0.363--0.386. Nó phải được revalidate nếu đổi
mass/model/propeller/plugin.

\subsection{Frame contract}

MAVLink \code{LOCAL\_NED} quy định north/east/down và z altitude âm
\cite{mavlinkCommon2026}. Các transform không được xem là đúng chỉ vì unit test
pass; convention và runtime data source cũng phải khớp.

\begin{table}[ht]
\centering\footnotesize
\caption{Interface contract cho coordinate frames.}
\begin{tabularx}{\textwidth}{L{0.20\textwidth}L{0.19\textwidth}L{0.23\textwidth}Y}
\toprule
\textbf{Source $\to$ destination} & \textbf{Convention/units} & \textbf{Code} & \textbf{Evidence/status}\\
\midrule
Gazebo world ENU $\to$ LOCAL NED & $(E,N,U)\to(N,E,-U)$, m và m/s & \code{enu\_to\_ned\_vel}, \code{ned\_to\_enu\_pos} & Spec + round-trip/sign test; live origin pending.\\
Body FLU $\to$ body FRD & $(x,y,z)\to(x,-y,-z)$ & point transform; \code{body\_flu\_rates\_to\_frd} & deterministic test.\\
Yaw NED $\to$ ENU & $\psi_E=\pi/2-\psi_N$ & \code{yaw\_ned\_to\_enu} & cardinal test.\\
Yaw-rate ENU $\to$ NED & $\dot\psi_N=-\dot\psi_E$ & \code{yaw\_enu\_to\_ned\_rate} & sign test.\\
Gazebo quaternion $\to$ ENU rotation & message xyzw; internal wxyz, Hamilton & \code{quat\_to\_rot} & identity/norm tests.\\
LiDAR/camera $\to$ body/world & SDF offset; odom full attitude & \code{scan\_to\_world\_enu} & source review; live frame pending.\\
\bottomrule
\end{tabularx}
\end{table}

\code{Odometry.twist.angular} hiện được giả định là body FLU. Exact installed
Gazebo plugin convention và một maneuver có dấu vẫn phải xác nhận trước rigid
live gate; đây là \textbf{UNVERIFIED/PENDING}.

\section{Tích hợp MAVLink và các điều kiện an toàn}

\subsection{Velocity target}

Vanilla và PA-MPPI v0 gửi \code{SET\_POSITION\_TARGET\_LOCAL\_NED}. Type mask
1479 bỏ position, acceleration và yaw; velocity cùng yaw-rate vẫn được sử dụng.
Mask đã được đối chiếu với enum của dialect \code{ardupilotmega} và handler
ArduCopter target commit \cite{mavlinkCommon2026,ardupilotSourceF808}. Dạng
symbolic là $1+2+4+64+128+256+1024=1479$; velocity và yaw-rate ignore bits đều
clear. Fallback velocity-only thêm yaw-rate-ignore 2048, tạo mask 3527.

\subsection{Goal tương tác từ RViz}

Node có thể subscribe \code{geometry\_msgs/PoseStamped} từ tool
\code{2D Goal Pose} trên topic \code{/goal\_pose}. Goal chỉ được chấp nhận khi
frame là \code{odom}; frame khác bị từ chối vì implementation chưa tự áp dụng
TF. Mỗi click thay route đang chạy, xóa terminal/recovery latch, reset output
conditioner và MPPI warm start. Vì tool 2D thường phát $z=0$, project bỏ giá trị
$z$ trong message và giữ altitude hiện tại, hoặc dùng altitude ENU cố định từ
\code{--rviz-goal-altitude}. Đây là \textbf{PROJECT INTERFACE DECISION} nhằm
tránh một thao tác XY vô tình tạo lệnh hạ xuống đất.

Trong paper profile, mỗi click tạo reference thẳng từ state hiện tại tới goal.
Do đó giao diện này thay việc nhập \code{Fly To}/tọa độ thủ công, nhưng không
phải global planner: operator vẫn phải đặt waypoint trung gian trong dead-end,
U-wall hoặc path bị che khuất. Callback topic và logic thay route đã được kiểm
tra offline; chuyến bay Gazebo điều khiển hoàn toàn bằng click còn PENDING.

\subsection{Thrust và body-rate target}

Rigid branch gửi \code{SET\_ATTITUDE\_TARGET} với mask 128: attitude quaternion
bị bỏ qua, cả ba body rate và thrust được giữ lại. ArduCopter không chấp nhận
body-rate vector thiếu một phần. Mask được dựng từ
\code{ATTITUDE\_TARGET\_TYPEMASK\_ATTITUDE\_IGNORE}, không phải magic decimal
\cite{mavlinkCommon2026}.

Có discrepancy phải giữ công khai: trang ``Copter Commands in Guided Mode''
hiện nói ba body-rate fields không được hỗ trợ và mask nên bằng 7
\cite{ardupilotGuidedDocs2026}. Ngược lại, source target commit
\code{f808f78...}, \code{GCS\_MAVLink\_Copter.cpp:892--966}, nhận quaternion,
nhận đủ cả ba rates, hoặc cho phép cả ba cùng ignored; partial rate vector làm
handler gọi hold \cite{ardupilotSourceF808}. Vì project khóa commit này,
implementation theo source behavior nhưng không che giấu khác biệt tài liệu.

Trường \code{thrust} chỉ mang nghĩa normalized thrust khi bit 3 của tham số
\code{GUID\_OPTIONS} được bật. Vì vậy cấu hình phải đặt giá trị này bằng 8.
Nếu giữ giá trị 0, ArduPilot diễn giải cùng trường thành climb-rate. Node đọc
lại parameter qua MAVLink và từ chối phát lệnh nếu giá trị không bằng 8. Bằng
chứng source là \code{mode.h:1208--1217}, \code{mode\_guided.cpp:651--655} và
handler lines 943--965 tại target commit \cite{ardupilotSourceF808}.

Các gate bổ sung gồm:

\begin{itemize}
  \item bắt buộc cờ \code{--experimental-attitude-control};
  \item chỉ hỗ trợ Gazebo odometry ở giai đoạn hiện tại;
  \item control frequency tối thiểu 20 Hz, khuyến nghị 50 Hz;
  \item thrust được giới hạn trong $[0,1]$ trước khi gửi;
  \item NaN/Inf trong rate hoặc thrust làm packet bị từ chối;
  \item planner vượt timeout cấu hình thì không phát control mới;
  \item stale, hard-brake, reached và Ctrl+C chuyển sang velocity-zero target,
  không phát zero thrust.
\end{itemize}

File \code{config/mppi\_attitude\_sitl.parm} đặt \code{GUID\_OPTIONS=8},
\code{MOT\_THST\_HOVER=0.38} và tắt hover learning cho thí nghiệm tái lập.

\section{Thiết kế kiểm thử}

\subsection{Automated tests}

Test suite dùng \code{unittest} và chạy trong env \code{ardupilot-rviz}. Các
nhóm test chính được liệt kê trong Bảng~\ref{tab:tests}.

\begin{table}[ht]
\centering
\caption{Phạm vi automated tests.}
\label{tab:tests}
\begin{tabularx}{\textwidth}{L{0.30\textwidth}Y}
\toprule
\textbf{Nhóm} & \textbf{Nội dung} \\
\midrule
Frame và MAVLink & ENU/NED, FLU/FRD, quaternion cơ sở, velocity mask 1479,
attitude mask 128, partial-rate rejection, GUID gate, NaN/Inf và saturation. \\
Planner safety & stale state, hard brake, timeout, reached, waypoint và diagnostics. \\
Occupancy/PA-MPPI & unknown/free/occupied, ray traversal, LOS, perception reward,
unknown collision và identical baseline dynamics. \\
Rigid-body & hover equilibrium, quaternion norm, body-rate response, thrust
mapping và một MPPI optimizer step. \\
Closed-loop clear-air & 80 chu kỳ tới goal 1 m, kiểm tra sai số đích và dải cao độ. \\
Configuration & YAML runtime fields và ưu tiên CLI override. \\
\bottomrule
\end{tabularx}
\end{table}

\subsection{Offline obstacle scenario}

Obstacle giả lập là một mặt stack container tại $x=12$ m, $y\in[-6,6]$ m và
$z\in[0,24]$ m. UAV bắt đầu tại $(0,0,20)$ m và bám tuyến waypoint
$(16,10,20)\rightarrow(30,0,20)$. Điều kiện đạt gồm:

\begin{equation}
\|\vect p_T-\vect p_{goal}\|<1\text{ m},
\qquad d_{min}\geq margin+1\text{ m}=5\text{ m}.
\end{equation}

Waypoint trung gian đóng vai trò global direction tối thiểu. Nó phá tính đối
xứng trái/phải trước obstacle rộng, một trường hợp mà local MPPI có thể tạo
trung bình control gần bằng không.

\section{Kết quả}

\begin{table}[ht]
\centering
\caption{Kết quả định lượng cập nhật ngày 13/09/2026.}
\label{tab:results}
\footnotesize
\renewcommand{\arraystretch}{1.08}
\begin{tabularx}{\textwidth}{L{0.31\textwidth}L{0.20\textwidth}Y}
\toprule
\textbf{Thí nghiệm} & \textbf{Kết quả} & \textbf{Chỉ số} \\
\midrule
Automated test suite & Đạt offline & 40/40 test; exact command lưu trong changelog. \\
Vanilla MPPI obstacle & Đạt offline & seed 7; final error 0.93 m; minimum clearance 6.30 m. \\
PA-MPPI v0 obstacle & Đạt offline & seed 7; final error 0.92 m; minimum clearance 6.11 m. \\
Global-path reference smoke & Đạt offline & seed 7; final error 0.23 m; minimum point clearance 5.85 m. \\
Paper-reference benchmark & Đạt offline & 3 map $\times$ seed 7/11/19: 9/9 terminal; point-radius analytical adapter. \\
Right-angle Gazebo/SITL & Đạt, một seed & 128 cycles; final error 0.235 m; path 23.15 m; min point clearance 2.19 m. \\
Narrow-gate Gazebo/SITL & Đạt, một seed & 192 cycles; final error 0.230 m; path 35.50 m; min point clearance 2.07 m. \\
Slalom Gazebo/SITL & Đạt, một seed & 193 cycles; final error 0.220 m; path 37.67 m; min point clearance 1.30 m. \\
Rigid hover dynamics & Đạt & position/velocity numerical error bằng 0;
quaternion norm bằng 1. \\
Rigid closed-loop unit gate & Đạt offline & 80 cycles; asserted final error $<0.5$ m và altitude trong 1.5--2.5 m; seed 7. \\
Rigid live Gazebo hover & Pending & Chưa chạy nhánh rigid body-rate/thrust trong Gazebo. \\
\bottomrule
\end{tabularx}
\end{table}

Ba run velocity-level live đều giữ cao độ trong khoảng 19.81--20.41 m và không
ghi deadline miss ở target 10 Hz. Đây chỉ là một seed trên mỗi map. Clearance
là khoảng cách tâm odometry--điểm LiDAR, chưa trừ footprint; slalom đạt 1.30 m,
thấp hơn collision radius rollout 1.5 m, nên chưa thể suy ra safety. Rigid
branch giữ hover trong internal model nhưng chưa có benchmark live/edge.

JSONL logger hiện ghi timestamp/state age, position, velocity, quaternion, body
rates, nominal/sent control, thrust, map status, clearance, goal/LOS, costs,
rollout statistics, selected trajectory, deadline/saturation/failsafe và seed.
\code{scripts/plot\_mppi\_experiment.py} chuẩn bị các plot trajectory, altitude,
velocity, commanded/measured rates, thrust, clearance, goal distance, latency và
MPPI-vs-PA comparison; không synthetic-fill field bị thiếu.

\section{Thảo luận}

\subsection{Mức độ đáp ứng yêu cầu mentor}

Kiến trúc hiện tại đáp ứng ý chính: obstacle avoidance và local planning nằm
ở companion computer; ArduPilot không cần obstacle ID và chỉ bám setpoint.
Vanilla MPPI là baseline chạy được. PA-MPPI v0 chứng minh occupancy/perception
cost. Rigid branch đã đưa control space tới thrust/body rate như hướng của bài
báo, nhưng mới ở mức offline-verified experimental path.

\subsection{Khoảng cách tới bản triển khai trong bài báo}

\begin{table}[ht]
\centering
\caption{So sánh phạm vi hiện tại với PA-MPPI \cite{zhai2026pamppi}.}
\footnotesize
\begin{tabularx}{\textwidth}{L{0.18\textwidth}L{0.29\textwidth}L{0.31\textwidth}Y}
\toprule
\textbf{Feature} & \textbf{Paper} & \textbf{Current implementation} & \textbf{Status/reason} \\
\midrule
State/control & $p,q,v,\omega$; thrust/rates & velocity $p,v,\psi$ or rigid $p,q,v,\omega$ & v0 reduced; rigid experimental.\\
Prediction & Full quad dynamics incl. torque/inertia & point-mass lag or thrust + first-order rate lag & Not full reproduction.\\
Map & ROG-Map; 3 states; 0.1 m, $5\times5\times2$ m & fixed NumPy/Torch grid; near-field heuristic & Backend differs.\\
Collision unknown & $G(p)\neq0$ penalized & same policy when map exists & Implemented offline.\\
Perception/axis & camera/body-x + goal ray & heading-x v0; body-x rigid & v0 approximation.\\
Ray/LOS phase & terminal DDA; two phases & terminal sampled ray; LOS switch & Different ray implementation.\\
Rollouts/horizon & 17,500 / 15 & 500/30 default; 128/20 ablation; 512/20 rigid & Reduced compute scale.\\
Timing & $\Delta t_{pred}=0.1$s, 50 Hz control & velocity 0.1 s/10 Hz; rigid 0.05 s/50 Hz target & Live deadline pending.\\
Framework & JAX + Agilicious & PyTorch + pytorch-mppi & Differs.\\
Mapping input & Depth + ROG-Map & Gazebo LiDAR hit endpoints & Sensor model differs.\\
Environment & Flightmare; C-wall/hole/4-wall; HIL/real & simple offline wall; Gazebo prepared & Matching evaluation pending.\\
Real flight & Demonstrated in paper & None & Pending hardware.\\
\bottomrule
\end{tabularx}
\end{table}

Do các khác biệt này, nhánh hiện tại phải được gọi là ``PA-MPPI v0'' hoặc
``rigid-body PA-MPPI experimental'', không phải reproduction đầy đủ.

\subsection{Rủi ro kỹ thuật còn lại}

\begin{itemize}
  \item Hover thrust phụ thuộc model, propeller và battery model. Giá trị 0.38
  phải hiệu chỉnh lại nếu SDF hoặc motor configuration thay đổi.
  \item Khối lượng 2.10 kg là tổng xấp xỉ từ model components; center of mass
  và inertia tổng chưa được nhận dạng chính xác.
  \item Gazebo odometry angular velocity cần được xác nhận thực nghiệm đúng
  body FLU trong mọi cấu hình plugin.
  \item Nhánh MAVLink state hiện chỉ đủ cho velocity planner; rigid branch chưa
  chuyển quaternion/rate từ estimator ArduPilot sang frame nội bộ.
  \item Sensor dropout, latency, obstacle động và estimator noise chưa được đưa
  vào Monte Carlo evaluation.
  \item Velocity-zero fallback đổi GUIDED submode nhưng không thay thế một
  independent safety controller hoặc người giám sát khi thử attitude control.
\end{itemize}

\section{Lộ trình kiểm chứng tiếp theo}

\begin{enumerate}
  \item \textbf{Hover gate:} takeoff bằng velocity/position controller, bật
  rigid node với goal ngắn trong không gian trống, theo dõi altitude, attitude,
  thrust target và mode.
  \item \textbf{Body-rate gate:} áp goal lệch 1 m theo từng trục, giới hạn rate,
  đo tracking delay để cập nhật $\tau_\omega$.
  \item \textbf{Clear-air route:} bay nhiều waypoint không obstacle; đo final
  error, overshoot, thrust saturation và deadline miss.
  \item \textbf{Static obstacle:} bật occupancy/perception cost với một stack
  container, sau đó tăng độ phức tạp môi trường.
  \item \textbf{Ablation:} so sánh vanilla, PA-MPPI v0 và rigid PA-MPPI bằng
  cùng seed/scenario theo success rate, clearance, path length, smoothness và
  compute p95.
  \item \textbf{Edge computer:} tích hợp timestamp và full attitude transform;
  sau đó benchmark tài nguyên và nhiệt độ.
  \item \textbf{Hardware gate:} HIL/tether test trước; chỉ flight test khi có
  geofence, mode switch và pilot override đã kiểm chứng.
\end{enumerate}

Simple wall có waypoint trung gian chỉ là baseline regression, không chứng minh
lợi ích perception. Fair experiment tương lai dùng C-wall/U-shaped occlusion với
cùng start, goal, world, dynamics, limits, horizon, base cost và seed; chỉ bật/tắt
perception terms. Hai config cố định là \code{config/ablation\_mppi\_unknown.yaml}
và \code{config/ablation\_pa\_mppi\_unknown.yaml}. Metrics bắt buộc gồm success
rate, time-to-goal, path length, minimum clearance, control smoothness, planner
cycles, compute mean/p95/worst, deadline misses, mapped free space và time tới
lần đầu goal visible. Toàn bộ gate/abort/log name nằm trong
\code{docs/GAZEBO\_RUN\_CHECKLIST.md}; hiện chưa gate nào được đánh dấu pass.

\section{Kết luận}

Repository đã chuyển được obstacle avoidance khỏi ArduPilot và hình thành một
pipeline local planning có thể thay đổi model/control level mà không sửa tầng
perception hoặc mission interface. Vanilla MPPI và PA-MPPI v0 đã chạy qua
obstacle scenario offline. Rigid-body experimental đã bổ sung quaternion,
thrust/body-rate control, \code{SET\_ATTITUDE\_TARGET}, calibration và safety
gate, đồng thời vượt qua 23 automated tests/offline gates. Kết quả đủ để chuyển
sang hover gate trong Gazebo, nhưng chưa đủ để tuyên
bố hoàn thành rigid PA-MPPI hoặc triển khai trên quadrotor thật.

\appendix
\section{Lệnh tái lập}

\subsection{Automated tests}

\begin{verbatim}
KMP_DUPLICATE_LIB_OK=TRUE OMP_NUM_THREADS=1 \
  /opt/miniconda3/envs/ardupilot-rviz/bin/python \
  -m unittest discover -s tests -p 'test_mppi*.py' -v
\end{verbatim}

\subsection{Rigid-body offline gate}

\begin{verbatim}
KMP_DUPLICATE_LIB_OK=TRUE OMP_NUM_THREADS=1 \
  /opt/miniconda3/envs/ardupilot-rviz/bin/python \
  scripts/mppi_velocity_avoidance.py \
  --planner rigid-pa-mppi \
  --config mppi_ardupilot/rigid_pa_mppi_config.yaml \
  --sim-test
\end{verbatim}

\subsection{Rigid-body live gate trong Gazebo}

\begin{verbatim}
MAVLINK20=1 KMP_DUPLICATE_LIB_OK=TRUE OMP_NUM_THREADS=1 \
  /opt/miniconda3/envs/ardupilot-rviz/bin/python \
  scripts/mppi_velocity_avoidance.py \
  --planner rigid-pa-mppi \
  --config mppi_ardupilot/rigid_pa_mppi_config.yaml \
  --experimental-attitude-control \
  --state-source odom --goal "1,0,20" \
  --mav tcp:127.0.0.1:5762 \
  --diag-jsonl output/log/rigid_hover_gate.jsonl
\end{verbatim}

\bibliographystyle{unsrt}
\bibliography{reports/pa_mppi_sources}

\end{document}
